\documentclass[11pt]{amsart}
\usepackage[style=alphabetic,backend=bibtex]{biblatex}
\addbibresource{bib.bib}

\input{preamble}
\usepackage{combelow}
\definecolor{electricindigo}{rgb}{0.44, 0.0, 1.0}
\definecolor{spirodiscoball}{rgb}{0.06, 0.65, 0.79}
\definecolor{warmth}{rgb}{0.89, 0.18, 0.42}
\definecolor{greenish}{rgb}{0.19, 0.78, 0.32}

\newcommand{\triv}{\mathrm{triv}}
\newcommand{\CKR}{C_{\mathrm{KR}}}
\newcommand{\HKR}{H_{\mathrm{KR}}}
\newcommand{\HY}{\YS \HKR}
\renewcommand{\CY}{\YS \CKR}
\newcommand{\ZU}{\EuScript{PS}}
\newcommand{\BSBim}{\mathrm{BSBim}}
\newcommand{\pTr}{\operatorname{pTr}}
\newcommand{\JM}{\operatorname{JM}}
\newcommand{\PS}{\EuScript{P}}
\newcommand{\xx}{\mathbf{x}}
\newcommand{\yy}{\mathbf{y}}
\newcommand{\Procesi}{\PS}
\newcommand{\Perf}{\mathrm{Perf}}
\newcommand{\bra}[1]{\langle #1 |}
\newcommand{\ket}[1]{|#1\rangle}
\newcommand{\ttheta}{\boldsymbol{\theta}}
\newcommand{\gl}{\mathfrak{gl}}
\newcommand{\BGGU}{\mathbf{BGG}_{\mathrm{univ}}}
\newcommand{\der}{\mathrm{der}}

\usetikzlibrary{scopes}

\begin{document}
%
\thispagestyle{empty}
\noindent{\Large Matthew T.~Hogancamp \hfill Research Statement}
\vskip5pt
\noindent

\noindent
%Northeastern University \hfill  email: m.hogancamp@northeastern.edu\\
%Department of  Mathematics  \hfill  office: Lake 447\\
%360 Huntington Ave, Boston, MA 02115  \vskip 4pt

%University of Southern California \hfill  email: hogancam@usc.edu\\
%Department of  Mathematics  \hfill  phone: (213) 740-3764\\
%3620 S. Vermont Ave., KAP 104 \hfill website: \href{http://www-bcf.usc.edu/~hogancam}{www-bcf.usc.edu/$\sim$hogancam}\\
%Los Angeles, CA 90089-2532\vskip 4pt


\hrule
\vskip14pt

My research interests are categorification, low dimensional topology, and representation theory.  I am particularly interested in homology theories in low dimensional topology, and their categorical / representation theoretic foundations.  I am driven by the rich structures afforded by link homology theories such as Khovanov homology and HOMFLY-PT homology, and their deep connections with geometric representation theory and algebraic geometry.



\subsection{Graphical calculus and knot Floer homology}
\label{s:OS}
In \cite{WebsterKIRT}, Webster defined a family link homology theories, associated to an arbitary complex semisimple Lie algebra $\mathfrak{g}$.  Webster's homology is isomorphic to the previously constructed Khovanov homology, respectively Khovanov--Rozansky homology, when $\mathfrak{g}=\gl_2$ (respectively $\mathfrak{g}=\gl_N$).

It is conjectured \cite{DGR} that there should also be link homology theories associated to certain Lie superalgebras, for instance Knot Floer homology should be $\gl_{1,1}$ link homology.  Unfortunately, it Webster's theory does not apply to Lie superalgebras, and it remains an open problem to construct $\gl_{N,M}$ link homology theories for $M,N\geq 1$ (the case $M=0$ or $N=0$ being accomplished already).  On the other hand, Ozsvath--Szabo have obtained a purely algebraic description of knot Floer homology, in terms of their so-called ``bordered knot algebras'' \cite{OSknot1,OSknot2}.  

During a reading course on bordered knot algebras, with my graduate student Connor Gallagher, it was realized that we could give a Webster-like flavor to these bordered knot algebras, and doing so revealed an action of Soergel bimodules that was not noticed before.

The main player is the following monoidal category that we introduce.
\newcommand{\reds}{\mathfrak{r}}
\newcommand{\blacks}{\mathfrak{b}}
\begin{definition}
Let $\CS$ be the $\Z$-graded monoidal category generated by objects $\reds$ (\emph{red}) and $\blacks$ (\emph{black}), with generating morphisms (and relations)
\[
R:=\begin{tikzpicture}[scale=.6,anchorbase]
\coordinate (a) at (0,0) ;
\coordinate (b) at (1,0) ;
\coordinate (c) at (0,1) ;
\coordinate (d) at (1,1) ;
\draw[red,very thick]
(a) -- (d);
\draw[black]
(b) -- (c);
\end{tikzpicture}
\qquad
L:=\begin{tikzpicture}[scale=.6,anchorbase]
\coordinate (a) at (0,0) ;
\coordinate (b) at (1,0) ;
\coordinate (c) at (0,1) ;
\coordinate (d) at (1,1) ;
\draw[red,very thick]
(b) -- (c);
\draw[black]
(a) -- (d);
\end{tikzpicture}
\qquad
u:=\begin{tikzpicture}[anchorbase]
\coordinate (a) at (0,0) ;
\coordinate (b) at (0,1) ;
\draw[red,very thick]
(a) -- (b);
\filldraw[red,very thick,fill=red]
(0,.5) circle (.2 em);
\end{tikzpicture}
\qquad \left(\text{modulo}\,
\begin{tikzpicture}[scale=.6,anchorbase]
\coordinate (a) at (0,0) ;
\coordinate (b) at (1,0) ;
\coordinate (c) at (0,1) ;
\coordinate (d) at (1,1) ;
\draw[black]
(b) -- (d)
(a) -- (c);
\end{tikzpicture} \ = \ 0
\qquad
\begin{tikzpicture}[scale=.6,anchorbase]
\coordinate (a) at (0,0) ;
\coordinate (b) at (1,0) ;
\coordinate (c) at (0,.75) ;
\coordinate (d) at (1,.75) ;
\coordinate (e) at (0,1.5) ;
\coordinate (f) at (1,1.5) ;
\draw[red,very thick]
(a) to[out=90,in=-90] (d) to[out=90,in=-90] (e);
\draw[black]
(b) to[out=90,in=-90] (c) to[out=90,in=-90] (f);
\end{tikzpicture}
=
\begin{tikzpicture}[scale=.6,anchorbase]
\coordinate (a) at (0,0) ;
\coordinate (b) at (1,0) ;
\coordinate (e) at (0,1.5) ;
\coordinate (f) at (1,1.5) ;
\draw[red,very thick]
(a) to (e);
\draw[black]
(b) to (f);
\filldraw[red,very thick,fill=red]
(0,.75) circle (.2 em);
\end{tikzpicture}\\
\qquad
\begin{tikzpicture}[scale=.6,anchorbase]
\coordinate (a) at (0,0) ;
\coordinate (b) at (-1,0) ;
\coordinate (c) at (0,.75) ;
\coordinate (d) at (-1,.75) ;
\coordinate (e) at (0,1.5) ;
\coordinate (f) at (-1,1.5) ;
\draw[red,very thick]
(a) to[out=90,in=-90] (d) to[out=90,in=-90] (e);
\draw[black]
(b) to[out=90,in=-90] (c) to[out=90,in=-90] (f);
\end{tikzpicture}
\ = \ 
\begin{tikzpicture}[scale=.6,anchorbase]
\coordinate (a) at (0,0) ;
\coordinate (b) at (-1,0) ;
\coordinate (e) at (0,1.5) ;
\coordinate (f) at (-1,1.5) ;
\draw[red,very thick]
(a) to (e);
\draw[black]
(b) to (f);
\filldraw[red,very thick,fill=red]
(0,.75) circle (.2 em);
\end{tikzpicture}\,
\right)
\]
Here we are employing the usual diagrammatics for discussing morphisms in monoidal categories, reading from bottom to top, so that for instance $R$ maps $\reds\blacks\to \blacks\reds$.
\end{definition}

\begin{remark}
The above graphical calulus looks like a hybrid of the ``redotted Webster algebras'' from \cite{KLSY}, and the homology of tangle Floer algebras from \cite{PetVer}.
\end{remark}


\begin{theorem}[In preparation, joint with grad student Connor Gallagher]\label{thm:OS Stendhal}
Let $\CS(m,k)$ be the full subcategory of $\CS$ consisting of objects with exactly $m$ occurences of $\reds$ and $k$ occurences of $\blacks$.  Let $A_{SOS}(m,k)$ denote the endomorphism algebra of the additive generator of $\CS(m,k)$.  Then there is an isomorphism $\mathcal{B}(m,k,\emptyset)\cong A_{SOS}(m,m+1-k)$, where $\mathcal{B}(m,k,\emptyset)$ is the algebra defined by Oszvath-Szabo in \cite{OSknot1}.
\end{theorem}


\begin{theorem}[In preparation, joint with grad student Connor Gallagher]\label{thm:Rouquier and OS}
There is an action of the category $\SBim_n$ of Soergel bimodules on $\CS(n,k)$ such that the induced categorical braid group action given by Rouquier complexes is quasi-isomorphic to the action constructed in \cite{OSknot1}.
\end{theorem}

\begin{remark}
Theorem \ref{thm:Rouquier and OS} is useful because it realizes the action of the braid group on $\CS(m,k)$ by dg functors, rather than $A_\infty$-functors.
\end{remark}

There are a lot of future directions to pursue here, and at least one of them  of them will surely be the subject of Connor's Ph.D.~ research.  Of these, the following is quite important.

\begin{problem}\label{prob:HHH to HFK}
Further elucidate the relationship between Soergel bimodules and knot Floer homology, obtaining a spectral sequence from triply-graded Khovanov--Rozansky homology to knot Floer homology. This would prove a conjecture of Dunfield--Gukov--Rasmussen \cite{DGR}.
\end{problem}

%\begin{problem}
%Extend the action of Soergel bimodules from Theorem \ref{thm:Rouquier and OS} to an action of singular Soergel bimodules in the style of the categorical symmetric Howe duality considered in \cite{KLSY}.  Deduce from this a colored version of knot Floer homology.
%\end{problem}



\subsection{Soergel bimodules, Hilbert schemes, and representabily of link homology}
\label{s:traces}
Another project of mine with concerns the construction of objects which represent link homology theories, much like how Eilenberg--Maclane spaces represent cohomology groups of spaces, in classical algebraic topology.  These representing objects ought to reveal some higher structure, in the same way that the Steenrod algebra of stable cohomology operations appears as the endomorphism algebra of the relevant Eilberg--Maclane spectrum.

For the purposes of this research statement I am concerned only with link homology theories in ``type $A$'', and the representing objects we seek will be certain complexes of Soergel bimodules.  Instead of Steenrod algebras, in this context we expect to find Rasmussen differentials and deep connections with the Hilbert scheme of points in $\C^2$.% (in accordance with a conjecture of Gorsky--Negu\cb{t}--Rasmussen).

Since link homology functors are ``trace-like'' we start with an investigation of the object which represents the \emph{universal} trace functor. 

In work in progress I show how to construct a complex $\ZU\in \Ch^b(\SBim_n)$ that I call the \emph{Procesi--Springer} complex because it places a role analogous to the Springer sheaf in character sheaves, and the Procesi bundle on the Hilbert scheme.  This complex satisfies a number of remarkable properties.

\begin{theorem}[In preparation]
We have $\Hom(\ZU,\one)\simeq \HH_\bullet(\SBim_n)$ and $\Hom(\one,\ZU)\simeq \HH^\bullet(\SBim_n)$.
\end{theorem}

It is known that $\HH_\bullet(\SBim_n)$ is formal \cite{HoLiformality} and isomorphic to the wreath product algebra $\k[\xx,\ttheta]\#S_n$ where $\xx=\{x_1,\ldots,x_n\}$ are even (polynomial) variables and $\ttheta=\{\theta_1,\ldots,\theta_n\}$ are odd (exterior) variables \cite{HRWtrace}.  On the other hand, computing Hochschild cohomology $\HH^\bullet(\SBim_n)$ is still an open problem (essentially equivalent to the problem of computing $\HH^\bullet$ of Soergel's graded version of BGG category $O(\mathfrak{\gl_n})$).

In \cite{GHMNserre}, my coathors and I proved that $\Ch^b(\SBim_n)$ admits a Serre functor, given by tensoring with $\FT\inv$, where $\FT$ is the Rouquier complex of the full twist braid.  In terms of $\ZU$ this has the following manifestation.


\begin{theorem}[In preparation]
The dual of $\ZU$ satisfies $\ZU^\vee\simeq \FT\inv\star \ZU$, in particular $\Hom(\ZU,\FT)\simeq \HH^\bullet(\SBim_n)$.
\end{theorem}

This should be compared to the familiar fact that the Procesi bundle $\PS$ on $\Hilb_n(\C^2)$ satisfies $\PS^\vee\cong O(-1)\otimes \PS$.  The following says that $\ZU$ represents the universal trace functor.

\begin{theorem}[In preparation]
The complex $\ZU$ has the structure of a coalgebra object in $\Ch^b(\SBim_n)$, as well has an object in the Drinfeld center of $\Ch^b(\SBim_n)$, so that $\Hom(\ZU,-)$ induces a functor from the horizontal trace $\hTr(\SBim_m)$ to the derived category of modules over $\Hom(\ZU,\one)\simeq \k[\xx,\ttheta]\#S_n$.  This functor is quasi-fully faithful, hence is quasi-equivalence onto  its image.  This image is the perfect derived category of $\k[\xx,\ttheta]\#S_n$.
\end{theorem}

\begin{remark}
The relation between $\hTr(\SBim_n)$ and modules over the wreath product algebra was known already \cite{HoLiformality,HRWtrace}, but the statements about $\ZU$ are new.
\end{remark}

There are a great many applications of $\ZU$ that I plan to investigate in future work.

\begin{problem}
Using the $S_n$ action on $\ZU$, for each $0\leq k\leq n$ define a direct summand $E_k\subset \ZU$ to be the $\sgn\times \triv$ isotypic component with respect to the restricted $S_k\times S_{n-k}$ action.  Show that $\Hom(E_k,-)$ is isomorphic to the $a$-degree $k$ part of triply graded Khovanov--Rozansky homology.
\end{problem}

\begin{remark}
In a parallel project, joint with Ben Elias, we are pursuing much more explicit formulae for $E_{k}$.  So far, we have succeeded in giving an explicit description $E_{1}$ using some new dg diagrammatics for Rouquier complexes.
\end{remark}

\begin{conjecture}
The object $E_\bullet=\bigoplus_{k=0}^n E_k$ has an action of an exterior algebra in variables $d_{N,M}$ for $N,M\geq 0$, so that adding $d_{N,M}$ to the differential of  $E_\bullet$ results in the object which represents $\gl_{N,M}$ link homology.
\end{conjecture}
\subsection{Invariants of manifolds coming from link homology}
\label{s:cat TV}

\newcommand{\BNs}{\mathrm{BN}}
\renewcommand{\BN}{\EuScript{BN}}


Perhaps the biggest open problem in the subject of link homology / categorification is to construct categorified TQFT-like invariants of manifolds in dimensions $\leq 4$.

There is a construction, essentially due to Bar-Natan which associates a vector space $\BNs(Y,L)$ to each pair $(Y,L)$ where $Y$ is a compact oriented 3-manifold and $L\subset \partial Y$ is a closed 1-submanifold of the boundary. In slightly more detail, $\BNs(Y,L)$ is formally spanned by embedded surfaces in $Y$ with boundary $L$, modulo some local relations.


Now, if $\Sigma$ is an oriented surface and $\pp\subset \partial \Sigma$ is a finite set of points, then we have the \emph{Bar-Natan category} $\BN(\Sigma,\pp)$; objects of this category are given by embedded 1-submanifolds $T\subset\Sigma$ with $\int T \subset \int \Sigma$ and $\partial T=\pp$, and $\Hom_{\BN(\Sigma,\pp)}(S,T)$ defined to be the $\BNs$-module associated to the thickened surface $\Sigma\times I$ with boundary condition $T\times\{1\}\cup B\times I\cup S\times\{0\}$.

%In an ongoing joint project with David Rose and Paul Wedrich, we are constructing a derived version of the Bar-Natan theory, and the resulting 3-manifold invariants give a categorical analogue of the Turaev-Viro quantum invariant for $\gl_2$ (with the quirk that we do not categorify roots of unity, so the decategorified construction is ill-defined).  In future work we plan to bump up the ambient dimension by 1, obtaining a derived, or chain level, version of the skein lasagna construction.


\renewcommand{\int}{\operatorname{int}}

Fix a commutative ring $\k$.  Let $Y$ be an oriented 3-manifold and $L\subset \partial Y$ a closed 1-submanifold of the boundary. Associated to the pair $(Y,L)$ we have the \emph{Bar-Natan module} $\BNs(Y,L)$, which is defined to be the $\k$-module spanned by isotopy classes of surfaces embedded in $Y$ with boundary $L$, modulo some local relations.


Now, if $\Sigma$ is an oriented surface and $\pp\subset \partial \Sigma$ is a finite set of points, then we have the \emph{Bar-Natan category} $\BN(\Sigma,\pp)$; objects of this category are given by embedded 1-submanifolds $T\subset\Sigma$ with $\int T \subset \int \Sigma$ and $\partial T=\pp$, and $\Hom_{\BN(\Sigma,\pp)}(S,T)$ defined to be the $\BNs$-module associated to the thickened surface $\Sigma\times I$ with boundary condition $T\times\{1\}\cup B\times I\cup S\times\{0\}$.

In \cite{} Rose, Wedrich and I show how to construct a dg category $\EuScript{FS}(\Sigma,\pp)$ which should be thought of as an appropriately derived version of the Bar-Natan category $\BN(\Sigma,\pp)$.  In work in progress we are working on the the following problem.

\newcommand{\SkLas}{\mathrm{SkLas}}
\newcommand{\SkLa}{\EuScript{S}k\EuScript{L}as}

%In joint work with David Rose and Paul Wedrich, we 

%\newcommand{\pp}{\mathbf{p}}
\newcommand{\FS}{\EuScript{F}}



%
%\begin{construction}
%To each compact oriented surface $\Sigma$ we have a dg category $\FS(\Sigma,\pp)$, depending on choice of a finite set $\pp\subset \partial \Sigma$ (possibly empty).
%%Let $\Sigma_1,\Sigma_2$ be closed oriented surfaces and $Y\colon \Sigma_1\to \Sigma_2$ an oriented cobordism.  Then 
%\end{construction}


\begin{problem}\label{prob:TV}
Construct a dg (or $A_\infty$) functor $\FS(\Sigma',\pp')\to \FS(\Sigma,\pp)$ associated to an oriented cobordism (with corners) $Y\colon \Sigma'\to \Sigma$.
\end{problem}

To achieve this, it would suffice to construct functors associated to 3-dimensional 0,1,2,3-handles and check the necessary relations.   We have well-defined functors for handles of index 0,1,2, and an idea for how 3-handles should work (with details still being developed).  

By comparing the construction with the Justin Roberts skein theoretic description of Turaev--Viro, we have the following. 

\begin{expectation}
The 3-manifold invariant constructed in a solution to Problem \ref{prob:TV} is a categorical analogue of the Turaev--Viro quantum invariant.
\end{expectation}






\printbibliography
\end{document}

\subsection{More Hilbert scheme predictions}
\label{ss:GNR}

The work of many authors \cite{} suggests a relationship between link homology (at least, those in type $A$) and various objects in algebraic geometry.  One of the most well known of these conjectures asserts that there ought to be a functor from the (derived) horizontal trace of $\SBim_n$ to the derived category of sheaves on the Hilbert scheme of $n$ points in $\C^2$, satisfying some list of expected properties.


To state the relevant conjectures and extract some concrete problems to focus on, we need a bit of background.


A closed point in $\Hilb^n(\C^2)$ is an ideal $I\subset \C[x,y]$ with $\dim(\C[x,y]/I)=n$.  There is an obvious action of $\C^\times \times \C^\times$, which acts by rescaling $x$ and $y$.  We will write $D^b(\Hilb^n(\C^2))$ for the derived category of $\C^\times \times \C^\times$-equivariant coherent sheaves. 

$\Hilb^n(\C^2)$ comes equipped with a rank $n$ vector bundle $T$, whose fiber at a point $I$ is the vector space $\C[x,y]/I$.  This vector bundle has a canonical section which (in the fiber at $I$) picks out $1\in \C[x,y]/I$, as well as a pair of commuting endomorphisms $X,Y$, given by the multiplication by $x$ and $y$. 

Secondly, the Hilbert scheme admits a very special vector $\Procesi_n$ of rank $n!$, called the \emph{Procesi bundle}.  This vector bundle satisfies
\begin{enumerate}
\item $\PS$ generates the derived category of $\Hilb^n(\C^2)$
\item the endomorphism algebra of $\Procesi_n$ is
\[
\End(\Procesi)\cong \k[x_1,\ldots,x_n,y_1,\ldots,y_n]\# S_n =: \k[\xx,\yy,S_n]
\]
and has no higher cohomology (so $\Procesi$ is a \emph{tilting generator}).
\end{enumerate}
These facts give an identification
\[
D^b(\Hilb^n(\C^2))\cong D^b(\text{mod-}\k[\xx,\yy]\#S_n).
\]

We have the following relationship between $\Procesi$ and $T$:
\begin{equation}\label{eq:P and T}
\Lambda^k(T) \cong (\PS)^{\sgn_k\times \triv_{n-k}}
\end{equation}
where we are restricting the $S_n$ action to $S_k\times S_{n-k}$ and taking the indicated isotopic component.


In particular the trivial component of $\Procesi$ is just $O$, and the sign component is the line bundle $O(1) = \Lambda^n(T)$.


\begin{conjecture}\label{conj:gnr 1}
There is a functor $\FS\colon \YS(\SBim_n)\to D^b(\Hilb^n(\C^2))$ such that the universal trace of $X$ is identified with $\Ext^\cdot(\PS,\FS(X)\otimes \PS)$ (as an object of $D^b(\k[\xx,\yy]\#S_n])$), and the triply graded homology of $X$ is identified with
\begin{equation}\label{eq:gnr conj}
\HH^\cdot(X) \simeq \Ext^\cdot(\Lambda^\cdot(T),\FS(Y)\otimes\PS)
\end{equation}
with the Hochschild grading on the left-hand side coinciding with the exterior algebra grading on the right-hand side.
\end{conjecture}

Note that a proof of the first part of this conjecture would be provided by the solution to Problem \ref{prob:ytrace} above.

Conjecture \ref{conj:gnr 1} inspires quite a few conjectures / problems that we discuss below.


\begin{problem}
Realize $E_{n,1}$ in terms of \emph{meridian operator} as sketched in the diagram below,
\[
dkjldkfj
\]
where we are invoking here a partial trace functor, which can be defined as the right adjoint to the inclusion $\SBim_n\to \SBim_{n+1}$.  This being done construct $E_{n,k}$ as the $k$-th exterior power of $E_{n,1}$.
\end{problem}

In a parallel project, joint with Ben Elias, we are pursuing much more explicit formulae for $E_{n,k}$.  So far, we have succeeded in giving an explicit description $E_{n,1}$ using some new dg diagrammatics for Rouquier complexes.  Finally, it is these $E_{n,k}$'s that should play a direct role in solution of Problem \ref{prob:ZnMN}.  Like much of the conjectures in this area, the following is motivated by Hilbert scheme wizardry.



\begin{conjecture}
The complex $E_{n,1}\in \Ch^b(\SBim_n)$ comes equipped with a pair of commuting endomorphisms $x_0,y_0$, and a map $\sigma\colon \one\to E_{n,1}$.  Furthermore, the representing objects $Z_{M,N}$ from Problem \ref{prob:ZnMN} can be constructed as
\[
Z_{M,N} = \left(E_{n,0}  \to E_{n,1} \to \cdots \to E_{n,n}\right)
\]
in which the differential is wedging with $x_0^My_0^N\sigma$ (using the  identification of $E_{n,k}$ with the $k$-th exterior power of $E_{n,1}$).
\end{conjecture}


%Given $C\in \YS_1(\SBim_n)$, let $C_0$ denote the underying complex of Soergel bimodules (whose differential is given by $\d_C|_{y=0}$).  Let $[C]:= [C_0]\otimes \k[y_1,\ldots,y_n]$ with differential given by $[\d_C] - \sum_i \theta_i y_i$.  \bigoplusLet $C\in \Ch^b(\SBim_n)$ be a complex of Soergel bimodules.  Suppose $C$ comes equipped with a deformed differential $\d_C+\a\in \End(C)\otimes\k[y_1,\ldots,y_n]$ with curvature $\d^2=\sum_i (x_i\star \id_C - \id_C\star x_i)y_i$.  L 


%
%
%\begin{problem}
%Construct an object $\PS\in \YS\SBim_n$ with the following structures / properties:
%\begin{enumerate}
%\item $\PS$ has the structure of an object in the Drinfeld center of $\YS\SBim_n$.
%\item $\PS$ is a coalgebra object.
%\item $\Hom(\PS,\one)\simeq \k[\xx,\yy,S_n]$.
%\end{enumerate}
%\end{problem}
%
%Note that such an object would determine a functor $\YS\SBim_n\to D^b(\k[\xx,\yy,S_n]\text{-mod})$ given by sending $X\mapsto \Hom(\PS,X)$, with the $\k[\xx,\yy,S_n]$-module structure given by the isomorphism $\k[\xx,\yy,S_n]\cong \Hom(\PS,\one)$ together with the coalgebra structure of $\PS$:
%\[
%\Hom(\PS,\one)\otimes \Hom(\PS,X) \to \Hom(\PS\star \PS, \one\star X) \to \Hom(\PS,X).
%\]
%
%The centrality of $\PS$ implies that $\Hom(\PS,-)$ is tracelike, in the sense that $\Hom(\PS,X\star Y)\simeq \Hom(\PS,Y\star X)$ naturally.
%
%\begin{conjecture}
%Assume Problem \ref{prob:P} is solved.  The functor $\Hom(\PS,-)\colon \YS\SBim_n \to D^b(\k[\xx,\yy,S_n])$ has essential image contained in $\Perf(\k[\xx,\yy,S_n])$, where as usual $\Perf$ denotes the dg category of perfect dg modules over a dg algebra.  Moreover, the resulting functor $\YS\SBim_n \to \Perf(\k[\xx,\yy,S_n])$ is the universal tracelike functor, meaning that it induces a quasi-equivalence between the derived horizontal trace of $\YS\SBim_n$ and $\Perf(\k[\xx,\yy,S_n])$.
%\end{conjecture}
%

%\begin{remark}
%Formal nonsense suggests that the object $\PS$ ought to be constructed as a deformation of the image of the trace of $\one_n$ under the left adjoint of the universal trace functor.  More on this later.
%\end{remark}
%
%
%\begin{problem}
%Construct objects which represent other traces on $\YS\SBim_n$ or $\SBim_n$ that we might care about, such as:
%\begin{enumerate}
%\item taking the Khovanov homology of the braid closure of $\b\in \Br_n$, which may be regarded as a tracelike functor from $\Ch^b(\SBim_n)$ to $\Ch^b(\k\text{mod}^\Z)$ (complexes of graded $\k$-modules)
%\item taking Batson--Seed homology \cite{} of the braid closure of $\b\in \Br_n$, which may be regarded as a tracelike functor from $\YS(\SBim_n)$.
%\end{enumerate}
%\end{problem}
%There are obvious generalizations of this, from Khovanov homology to $\gl(N)$-link homology (of which Khovanov homology is the $N=2$ special case).  It is conjectured that there should be link homology theories associated to the super Lie algebras $\gl(M|N)$ for all $M,N\in \Z_{\geq 0}$, with the $\gl(1|1)$-link homology being isomorphic to an appropriately ``unreduced'' version of link Floer homology (see \cite{}).  One approach to constructing such link homologies would be the construction of the objects which represent them:
%
%\begin{problem}\label{prob:ZnMN}
%For each $n,M,N\geq 0$ construct an object $Z_n(M|N)\in \YS\SBim_n$ so that $\Hom(Z_n(M|N),-)$ yields a well-defined link homology theory, so that the case of $Z_n(M|0)$ recovers the (appropriately deformed version of) $\gl_M$ Khovanov--Rozansky homology, and the case of $Z_n(1|1)$ recovers Oszvath--Szabo link homology from \cite{}.
%\end{problem}
%
%Or we could ask for something much more basic:
%
%\begin{problem}\label{prob:Enk}
%Given integers $0\leq k\leq n$, construct an object $E_{n,k}\in \YS\SBim_n$ which represents the Hochschild degree $k$ component of the deformed triply graded link homology $\HY$.
%\end{problem}
%
%
%\begin{conjecture}\label{conj:Enk and P}
%There is a map of dg algebras $\k[\xx,\yy,S_n]\to \End(\PS)$ (not an isomorphism.  Taking appropriate isotopic component with respect to the $S_n$ action yields the objects $E_{n,k}$ from Problem \ref{prob:Enk}.  Precisely we have
%\begin{equation}\label{eq:Enk and P}
%E_{n,k} \simeq \langle \PS, e_k h_{n-k}\rangle
%\end{equation}
%\end{conjecture}
%%In \eqref{eq:Enk and P} we are using the notation that if $\chi$ is the character of an $S_n$-representation $V$, $\langle{\PS,\chi}=\PS\otimes_{\k[S_n]} V$ denotes the associated isotopic component.   Furthermore $e_kh_{n-k}$ is the product of an elementary symmetric function $e_k$ and the a complete homogeneous symmetric function $h_{n-k}$; this product is the sum of two Schur functions associated to hook shapes.
%
%
%Conjecture \ref{conj:Enk and P} is motivated by connections to the Hilbert scheme of points in $\C^2$ (also conjectural).  We will discuss these in \S \ref{}.  The same connections to the Hilbert scheme inspire the following.
%
%\begin{conjecture}\label{conj:Enk and En1}
%The braid group action on $E_{n,1}^{\star k}$ coming from the Drinfeld central structure factors through the symmetric group, so that one can define Schur functors acting on $E_{n,1}$.  Furthermore, the $k$-th exterior power of $E_{n,1}$ is $E_{n,k}$:
%\[
%E_{n,k}\simeq \Lambda^k(E_{n,1}) := \langle E_{n,1}^{\star k}, e_k\rangle
%\]
%\end{conjecture}
%
% 
%The question of representability of a given link homology theory is solved completely formally.
%
%\begin{theorem}\label{thm:res of diag}
%The category $\SBim_n$ is smooth, in the sense that it admits a finite projective resolution as a bimodule over itself.  In fact, we have a projective resolution of $\SBim_n$ of the form
%\[
%\widetilde{\bigoplus_{w\in S_n}} \ket{\Delta_w}\otimes_R^L \bra{\nabla_w}
%\]
%where $\Delta_w$ (resp.~$\nabla_w$) is the Rouquier complex associated to a positive (resp.~negative) braid lift of a permutation $w\in S_n$, the symbol $\otimes_R^L$ denotes the derived tensor product over the ring $R=\k[\xx]$, and the symbols $\ket{\ }$, $\bra{\ }$ denote the left and right Yoneda modules:
%\[
%\ket{X} = \Hom(X,-) \ , \qquad \bra{Y} = \Hom(-,Y).
%\]
%\end{theorem}
%This theorem gives, in particular a BGG-style resolution of any complex $X\in \Ch^b(\SBim_n)$ of the form
%\[
%X\simeq \widetilde{\bigoplus_{w\in S_n}} \Hom(\Delta_w,X)\otimes_R \bra{\nabla_w}
%\]
%(where we have dropped the derived tensor product because homs between Soergel bimodules happen to be free over $R$).  Another very useful consequence of Theorem \ref{thm:res of diag} is that it gives the f=ollowing (this is basically Brown representability, but finiteness properties allow us to stay in the realm of bounded complexes).
%
%\begin{theorem}
%Any dg functor $F\colon \Ch^b(\SBim_n)\to \DS$ has a left (or right) adjoint, defined on objects by
%\[
%LF(Y)=\widetilde{\bigoplus_{w\in S_n}} \ket{\Delta_w}\otimes_R^L \Hom_\DS(Y,F(\nabla_w))^{\ast}
%\qquad\quad
%RF(Y)=\widetilde{\bigoplus_{w\in S_n}} \Hom_{\DS}(\Delta_w,Y)\otimes_R^L \bra{\nabla_w}
%\]
%(here $M^\ast=\Hom_R(M,R)$ denotes the linear dual of an $R$-module $M$).
%\end{theorem}
%
%\begin{corollary}
%Let $F$ be a link homology theory which factors through $\Ch^b(\SBim_n)$.  This link homology is represented by an object of the form
%\[
%Z_n(F) = \tilde{\bigoplus_{w\in S_n}} \Delta_w\otimes_R^L F(\nabla_w)^\ast
%\]
%\end{corollary}
%
%
%\MH{Todo:
%\begin{enumerate}
%\item Something about yified functors.
%\item What kind of progress have I made?  Well i can construct $\PS$, and I think the comultiplication is dual to the multiplication on $\PS^\vee$, which comes from the shuffle product on the bar complex, essentially.  I think I can show that $\PS$ is central, and that $\Hom(\PS,\one)$ is what it needs to be.
%\item With Eugene we can construct a complex that ``looks like'' it should be a model for $E_{n,1}$.  It has 
%\item With Ben we can construct $E_{n,1}$ explicitly.
%\end{enumerate}
%}
%
%
%
%\subsection{Connecting link homology with Hilbert schemes}
%
%A well-known family of conjectures \cite{GNR} relates the link homology $\HKR(L)$ with the Hilbert scheme $\Hilb^n(\C^2)$ of points in $\C^2$. In this section I will outline some work I have done in shedding light on this connection, and some future work I have planned.
%
%
%
%% The joint spectrum of $X$ and $Y$ determines a point in $\Sym^n(\C^2)$; the resulting map $\Hilb^n(\C^2)\to \Sym^n(\C^2)$, called the Hilbert--Chow map, is a resolution of singularities.  There is a closely related scheme called the \emph{isospectral Hilbert scheme $X_n$, which morally speaking is obtained from $\Hilb^n(X)$ by choosing an ordering on the multiset of joint eigenvalues.  Technically, $X_n$ is the reduced fiber product of $(\C^2)^n$ and $\Hilb^n(\C^2)$, over $\Sym^n(\C^2)$.  Haiman showed that $X_n$ is flat over $\Hilb^n(\C^2)$, so th
%
%
% Given a sheaf $\FS$ on $\Hilb^n(\C^2)$ we may consider its quadruply-graded cohomology defined by
%\begin{equation}\label{eq:quad homology}
%H^{\cdot\cdot\cdot\cdot}(\FS) := \Ext^\cdot(\Lambda^\cdot(T), \FS)^{\cdot,\cdot}
%\end{equation}
%where $\Lambda^\cdot(T) = \bigoplus_{k=0}^n\Lambda^k(T)$ denotes the exterior algebra.  The four gradings are the cohomological degree coming from $\Ext^\cdot$, the exterior algebra degree, and two gradings coming from the $\C^\times \times \C^\times$ action.
%
%
%Notice that the functor $H^{\cdot\cdot\cdot\cdot}(-)$ is representable, since it is homs out $\Lambda^\cdot(T)$ (in $D^b(\Hilb^n(\C^2))$).  As a first step in elucidating \eqref{eq:gnr conj}, we ask whether $\HY(\b)$ is representable in a similar way.
%
%
%\subsection{Representability of link homology}
%\label{ss:representability}
%
%
%
%\begin{problem}
%Construct a ``tracelike'' functor $\YS\SBim_n \to D^b(\k[\xx,\yy,S_n])$ and its left adjoint $ \YS\SBim_n\leftarrow D^b(\k[\xx,\yy,S_n])$, with essential image contained inside the Drinfeld center of $\YS\SBim_n$.
%\end{problem}
%
%\begin{problem}
%Construct a Drinfeld central object $\PS_n\in \YS\SBim_n$ with endomorphism algebra (quasi-isomorphic to) $\k[\xx,\yy,S_n]$, so that $\Hom(\PS^{\SBim}_n,-)$ is the desired functor $\YS\SBim_n\to D^b(\k[\xx,\yy,S_n])$, with left adjoint $M\mapsto M\otimes_{\k[\xx,\yy,S_n]} \PS^{\SBim}_n$.
%\end{problem}
%
%
%The functor
%
%\subsection{Representing link homology}
%
%Consider the symmetric group $S_n$. The transposition (or reflection) which swaps $i$ and $j$ will be denoted $(i,j)$.  The simple reflections are the the elements $s_i:=(i,i+1)$ for $1\leq i\leq n-1$.  The length of $w$ is the smallest number of simple reflections in an expression $w=s_{i_1}\cdots s_{i_l}$.  The reflection length of $w$ is the smallest number of (not necessarily simple) reflections in an expression $w=r_{i_1}\cdots r_{i_l}$.  The reflection length is invariant under conjugation, hence depends only on the cycle type of $w$.  In fact, the reflection length of $w$ equals $n-c$, where $c$ is the number of cycles in $w$.
%
%The Bruhat order on $S_n$ is the partial order generated by relations $w\leq sw$ if $\ell(w)<\ell(sw)$ (and $w\leq ws$ if $\ell(w)<\ell(ws)$) for all simple reflections $s$.
%
%\subsection{Rouquier complexes}
%
%
%\begin{construction}
%To each braid $\b\in \Br_n$ there is a curved Rouquier complex $\tw T(\b)\in \YS_{n,w}$.  To each permutation $w\in S_n$ there is a Kozul matrix factorization $K_w\in \YS_{n,w}$ (whose dimension equals the reflection length of $w$).  Then if $w\in S_n$ is the permutation represented by $\b$ we construct $\HY(K_{w\inv}\star \tw T(\b))$, whose homology depends only on the link $L=\hat{\b}$ up to isomorphism and shift in the trigrading. 
%\end{construction}
%
%\begin{definition}
%Splitting map.
%\end{definition}
%
%\begin{theorem}\label{thm:FT and Hilb}
%The map $\HY(\FT_n^k)\to \HY(\one_n)$ induced from the splitting map is injective and induces an isomorphism of graded algebras $\bigoplus_{k\geq 0}\HY(\FT_n^k)$ Moreover, the ima
%\end{theorem}
%
%\begin{expectation}
%\begin{itemize}
%\item The Procesi bundle corresponds to $\widetilde{\bigoplus}_{w\in S_n} \Delta_w\star \Delta_{w\inv}$.
%\item The ideal sheaf corresponds to $\pTr^1(J_{n+1})$.
%\item The tautological vector bundle corresponds to $\pTr^1(\Cone(J_{n+1}\to \one_{n+1}))$.
%\item The canonical line bundle $\mathcal{O}(1)$ corresponds to $\FT_n$. 
%\end{itemize}
%\end{expectation}
%
%\begin{theorem}
%There exists a complex $\mathbb{P}_n=\widetilde{\bigoplus}_{w\in S_n} \nabla_w \otimes_R \Lambda\otimes_R \nabla_w\in \Ch^b(\SBim_n\otimes_\k \SBim_n)$ where $\Lambda$ is the Koszul resolution of $R$ by free graded $R\otimes_\k R$-modules.  This complex satisfies:
%\begin{enumerate}
%\item idempotent
%\item a functor $F\colon \SBim_n\to \DS$ has a right adjoint sending $D\mapsto \widetilde{\bigoplus}_{w\in S_n} \Hom_{\DS}(F(\nabla_w^\vee),D)\otimes_R^L \nabla_w$.
%\item more?
%\end{enumerate}
%\end{theorem}
%
%\begin{construction}
%Let $E_{n,k}\in \Ch^b(\SBim_n)$ be the complex obtained by applying $\HH^k$ to the first tensor factor of $\mathbb{P}_n$, i.e.~
%\[
%E_{n,k}^\vee:= \widetilde{\bigoplus}_{w\in S_n} \HH^k(\nabla_w) \otimes_R \nabla_w\in \Ch^b(\SBim_n\otimes_\k \SBim_n)
%\]
%\end{construction}
%
%\[
%\begin{tikzpicture}[scale=.5]
%\coordinate (a0) at (0,0);
%\coordinate (a1) at (1,0);
%\coordinate (a2) at (2,0);
%\coordinate (a3) at (3,0);
%\coordinate (a4) at (4,0);
%\coordinate (a5) at (5,0);
%\coordinate (a6) at (6,0);
%\coordinate (b0) at (0,3);
%\coordinate (b1) at (1,3);
%\coordinate (b2) at (2,3);
%\coordinate (b3) at (3,3);
%\coordinate (b4) at (4,3);
%\coordinate (b5) at (5,3);
%\coordinate (b6) at (6,3);
%\coordinate (c) at (-1,1.5);
%\draw[thick]
%(a3) .. controls ++(0,1) and ++(0,-1)..  (c);
%\draw[line width=2.5mm,white]
%(a0) -- (b0)
%(a1) -- (b1)
%(a2) -- (b2)
%%(a3) -- (b3)
%(a4) -- (b4)
%(a5) -- (b5);
%%(a6) -- (b6)
%\draw[thick]
%(a0) -- (b0)
%(a1) -- (b1)
%(a2) -- (b2)
%%(a3) -- (b3)
%(a4) -- (b4)
%(a5) -- (b5);
%%(a6) -- (b6)
%\draw[line width=2.5mm,white]
%(c)  .. controls ++(0,1) and ++(0,-1).. (b3);
%\draw[thick]
%(c) .. controls ++(0,1) and ++(0,-1).. (b3);
%\node at (0,-.5) {\scriptsize{$1$}};
%\node at (1.5,-.5) {\scriptsize{$\cdots$}};
%\node at (3,-.5) {\scriptsize{$i$}};
%\node at (4,-.5) {\scriptsize{$\cdots$}};
%\node at (5,-.5) {\scriptsize{$n$}};
%\end{tikzpicture}
%\]
%
%\begin{problem}
%Construct finite complexes $E_{n,k}$ in the Drinfeld center of $\Ch^b(\SBim_n)$ (rather, the $A_\infty$ Drinfled center as in \cite{})
%\begin{enumerate}
%\item on the level of Grothendieck groups, $[E_{n,k}]=\sum_{1\leq i_1<\cdots<i_k\leq n} J_{i_1}\cdots J_{i_k}$ (the $k$-th elementary symmetric function evaluated on the Jucys-Murphy braids).
%\item the hom complex $\Hom_{\Ch^b(\SBim_n)}(E_{n,k},B)$ is homotopy equivalent to $\HH^k(B)$, naturally in $B$.
%\item \MH{more?}
%\end{enumerate}
%\end{problem}
%
%Note that the $E_{n,k}$ are uniquely characterized, up to homotopy equivalence, by property (2). The special cases $E_{n,0}=\one$, $E_{n,n}=\FT_n$ are easy. For $E_{n,1}$, we have the following.
%
%\begin{expectation}
%The representing complex $E_{n,1}$ can be constructed as
%\[
%E_{n,1} \simeq I^R(\Cone(J_{n+1}\to \one_{n+1}))
%\]
%where $I^R$ is the right adjoint to $I\colon \SBim_n\colon \SBim_{n+1}$ defined by $B\mapsto B\boxtimes \one_1$, and the map $J_{n+1}\to \one_{n+1}$ is the so-called \emph{splitting map} from \cite{}.
%\end{expectation}
%
%For $E_{n,k}$ in general we have the following.
%\begin{expectation}
%Tthe braid group action on $E_{n,1}^{\otimes k}$ (coming from the Drinfeld central structure) factors through the symmetric  group up to homotopy, and $E_{n,k}\simeq \Lambda^k(E_{n,1})$.
%\end{expectation}

%
%\bibliographystyle{alpha}
%\bibliography{bib}
